
% \documentclass[conference]{IEEEtran}
% \documentclass[journal,12pt,onecolumn,finalnofoot,]{IEEEtran}
\documentclass[journal,12pt,onecolumn,draftclsnofoot,]{IEEEtran}
\IEEEoverridecommandlockouts
% The preceding line is only needed to identify funding in the first footnote.
% If that is unneeded, please comment it out.
\usepackage{cite}
\usepackage{amsmath,amssymb,amsfonts}
\usepackage{algorithmic}
\usepackage{graphicx}
\usepackage{textcomp}
\usepackage{enumitem}
\def\BibTeX{{\rm B\kern-.05em{\sc i\kern-.025em b}\kern-.08em
    T\kern-.1667em\lower.7ex\hbox{E}\kern-.125emX}}
\begin{document}

\title{ Comparison and Evaluation of  state-of-art Blockchain Scaling
Solutions}

\author{\IEEEauthorblockN{Nakul Chawla} \\
\IEEEauthorblockA{\textit{Computer Science} \\
\textit{Arizona State University}\\
nchawla3@asu.edu}
}


% \listadd{\committeemember}{Member Name}


\maketitle

\centerline{Thesis Proposal} 
\centerline{for the Master's Degree in Computer Science}

\hfill \break
\hfill \break
\hfill \break
\hfill \break
\hfill \break
\hfill \break

\centerline{\textbf{Approved by:} }
\centerline{Selcuk K. Candan (Chair)} 
\centerline{Dragan Boscovic (Chair) }

\newpage


\begin{IEEEkeywords}
    Blockchain, Scalability, XThin Blocks, Compact Blocks,Graphene Blocks, Lightning Network, Plasma, Bitcoin,
    Ethereum, Dash
\end{IEEEkeywords}

\section{Introduction}
In 2009, Bitcoin\cite{bitcoin} introduced the first blockchain which is an
append only database that stores transactions. Blockchain maintains the
distributed database in a decentralized network aiming to solve the Byzantine
Generals Problem \cite{bgp}. Byzantine Generals Problem is a problem in
distributed computing in which two generals have to come to an agreement on
whether to attack or not, but have only message passing as a way of
communicating. This problem in Finance is extrapolated to digital money by
introducing malicious nodes into the system who would want to double spend their
money. In a typical banking system, the bank acts as the central node which
validates every transaction that is made by the user. However, in a peer-to-peer
system, where no central node validates the transaction, malicious/byzantine behaviour is
easy to perform. 
Bitcoin solves this problem by introducing a shared public ledger called
blockchain, in which all confirmed transactions are included.   

\subsection{Process}
Nodes on the bitcoin network arrange themselves in a decentralized manner
connecting to a anonymous peers to either previously known IP
addresses or by DNS seeds. There are several types of nodes: Miners, 
Full-Nodes and SPV (Simplified Payment Verification) wallets. Full-nodes and SPV
clients are required to propagate and verify the transactions, where as miners
are required to verify and add transactions to a block and add it to the
blockchain. Once valid a block propagates through the network, it is then added to the
blockchain by mutual consensus of the network.

\subsection{Problems}
Blockchain currently faces a number of limitations that are listed below:

\subsubsection{Scalability}
An arbitrary time interval of ten minutes was set between adding two blocks when
bitcoin was introduced. Along with a ten minute interval, a 1MB block size was
set later after a DoS attack in 2010. These parameters have proven to constrain
the throughput of bitcoin network to 3-7 transactions per second. There have
been several techniques to change this, but none have yet proved to scale it to
the current fiat currency vendors VISA that is approximately 24k transactions
per second during peak.

\subsubsection{Storage}
Over the time as the size of blockchain grows, individual contributors i.e. the
full nodes who volunteer verification of the block will face trouble validating
each transaction on the network.The current size of bitcoin blockchain in close
to 160GB, which all full nodes store on their personal storage devices. As
blockchains get widely accepted, and replaces the current credit/debit cards,
the estimated size of blockchain per year required will be ~400 terabytes as
stated in \cite{lightning}, which
individual volunteer nodes won't be able to afford. 


% \section{Background And Related Work}
%
% \subsection{Background}
%     Blockchain is a decentralized, tamper-proof digital ledger technology that
%     is transforming transaction prospects for many industries. In addition to
%     being the foundation of cryptocurrencies like Ethereum, the blockchain
%     structure, which consists of linked blocks of information, allows for direct
%     transactions across a network of computers without the need for a central
%     authority. The Ethereum network consists of three kinds of nodes:
%     \textbf{full nodes}, \textbf{light-weight nodes } and \textbf{miners}. Nodes
%     save the complete blockchain and have two main functions: 
%     \begin{enumerate}
%         \item to validate transactions
%         \item to propagate transactions and blocks through the network
%     \end{enumerate} 
%     Although full nodes validate transactions by checking their inputs and
%     outputs, light-weight nodes verify transactions by SPV, or by requesting random
%     chunks of a block and using the merkle root to verify each chunk. This
%     verification technique relies on trusting a full node, or whichever node is sending
%     data. This can be dangerous, because a malicious node can withhold small
%     amounts of data and can get unnoticed by many light clients. The network
%     becomes much more trustworthy if the number of full nodes are a majority
%     since data is always verified. As blockchain technology becomes increasingly
%     adopted, the storage requirements are also increasing for each individual
%     node. This
%     increase in infrastructure requirement forces volunteer nodes to move to
%     light weight nodes. This harms the network in two ways :
%     \begin{enumerate}
%         \item The network gets more centralized to people who can afford
%             supercomputers
%         \item The trust on the network weakens as most transactions are not
%             completely validated by full nodes
%     \end{enumerate}
%     In this project, we want to introduce a new kind of node that saves only a
%     part of the blockchain; thus, reducing the storage requirement while also
%     validating each transaction that is a part of the block. This will introduce
%     the following advantages to the network:
%     \begin{enumerate}
%         \item The trust on the network will be high as the nodes will validate
%             each transaction fully
%         \item Randomizing block storage in such a way that transactions are
%             fully validated by a small part of the network leading to multiple
%             transactions being validated in parallel.
%     \end{enumerate}
%     We will be focusing only on the first advantage as a part of this project,
%     the second will remain a future work.
%
% \subsection{Related Work}

\section{Proposed Solutions}
Below are some of the solutions that are being adopted slowly as potential solution.

\subsection{On-chain scaling}

\subsubsection{Increasing Block Size}
A very naive approach of solving the scalability problem is to directly increase
the block size to what is currently required to handle all transactions in the
mempool. The fallbacks of this plan however can affect the network security and
can end up being an economical disaster by following means:

\begin{enumerate}
    \item Increasing block size directly increases the percentage of orphan
        blocks because of the delay of propagating the block through the network
        in the given block interval.More orphan rates harm the economic
        incentive of miners. Discouraging new miners can eventually centralize
        the network.

    \item Increasing block size directly affects the size of blockchain that is
        stored by each node of the network. A substantial increase in the size
        of blockchain discourages volunteer full-nodes on the network and
        centralizes the propagation to parties that can afford supercomputers.

    \item Increasing block size linearly only increases the transaction count
        linearly. In order to match up to VISA like vendors, there needs to be
        an exponential rise in the number of transactions. 

\end{enumerate}

\subsubsection{Masternodes}
Dash \cite{dash} coin was introduced as Xcoin in the year 2014. It introduced the first
on-chain multi-tier scaling solution by introducing what they call Masternodes.
These masternodes are not volunteer nodes and require to have a higher
configuration (bandwidth and processing). In order to fulfill these
requirements, they are awarded 45\% of the mining reward as compensation. The masternodes add
reliability to the network, and because of the higher configuration push the
network to produce higher throughput.The block interval is also reduced to
two minutes and 30 seconds and recently increased the block size to 2MB. This
arrangement helps in packing double the transaction amount in a much lesser time
which increases the throughput 8 fold when the network will reach its peak.

\subsubsection{Segregated Witness \cite{segwit}}
Segregated Witness, also called SegWit, is a software update that tries to fix
transaction malleability which had been a known weakness in bitcoin
security.This method segregates the signature data from the original data. 
Although, this method wasn't directly aimed at solving scalability, it reduces the
size of each transaction, thus packing more transactions into the block. 

\subsubsection{Compact \cite{compact}, XThin \cite{xthin} and Graphene blocks
\cite{graphene} }
Since increasing block size has direct effect on centralization of the network,
a more novel approach is to utilize the fact that the nodes on the network
already have most of the information that is needed for the new block to be
formed. All that is needed is the metadata for the other nodes to form this new
block. \textbf{Compact blocks} instead of sending the complete block, send only
the short transaction IDs that are 6 bytes each in the block along with the
coinbase transaction. This compresses the block substantially from a 1MB to a few
kBs. Since the information that needs to be propagated is much smaller now, more
information can be packed in the same 1MB block, effectively increasing the
throughput of the network. \textbf{XThin blocks} use the same technique of
compression, but instead of sending shortIDs of the transactions, sends a bloom
filter of the complete mempool. From this bloom filter sketch, the second node
can easily know what information the node does not have and only send that
information. Since, bloom filters have a false positive rate that increases as
the amount of data increases in it, \textbf{Graphene blocks} along with a bloom
filter also send an inverted bloom lookup table (IBLT) in order to check the
false positive rate and correctly form the group of transactions that are
included in the block being propagated.

\subsection{Off-chain scaling}
\subsubsection{Bitcoin Lightning Network \cite{lightning}}
Bitcoin lightning network is based on the idea of bidirectional off-chain
payment channels. Two participants create a ledger entry on the blockchain which
requires both participants to sign off on any spending of funds.
\cite{lightning} explains the protocol as "Both parties
create transactions which refund the ledger entry to their individual
allocation, but do not broadcast them to the blockchain. They can update their
individual allocations for the ledger entry by creating many transactions
spending from the current ledger entry output. Only the most recent version is
valid, which is enforced by blockchain-parsable smart-contract scripting. This
entry can be closed out at any time by either party without any trust or
custodianship by broadcasting the most recent version to the blockchain."

Once these channels are made by multiple parties, a node can find an already
existing path to pay another node thus avoiding the need to go on chain and
explicitly adding transactions to a block. Since, the nodes on the path can't be
trusted, atomicity is introduced with decrementing Hashed time-locks.

This results in making payments off-chain without any constraints in a limitless
manner. In an event of discrepancy, any one party can broadcast the previous transaction
to the blockchain adding it to the immutable database that can't be changed. The
malicious node, however is penalized with all funds in the transaction to be
transferred to the other party, thus preventing byzantine failures.


\subsubsection{Ethereum Plasma \cite{plasma}} 
Ethereum\cite{ethY} \cite{ethW} currently faces problem of scalability with the number of users on its
platform. In order to scale to a size where Ethereum can handle API requests as
large 100k per second, the plasma whitepaper details a way to form side chains
called the plasma chains. Plasma is composed of two key parts:

\begin{enumerate}
    \item Reframing all blockchain computation into a set of MapReduce functions
    \item Optional method of Proof-of-Stake token bonding that realizes Nakamoto
        consensus which will penalize malicious actors. 
\end{enumerate}

A tree hierarchy of many blockchains will be formed, that treats each individual
branch as a blockchain that is enforceable with the main chain and MapReducable
computation committed into merkle proofs.

\section{Intent}
The thesis intention is to evaluate the solutions provided above by implementing
their respective simulator and emulator. The specific criteria for evaluating
each technique is listed below in the section Evaluation Criteria.

\subsection{Components}
The thesis will encompass two projects:

    \subsubsection{The Simulator}
    \begin{enumerate}[label=(\alph*)]
    \item 
            The bitcoin simulator is an NS3 based
            simulator that was originally built for the research shown in
            \cite{security}. The simulator demonstrates the working of bitcoin, dogecoin and
            litecoin, but bases it's research on block being the fundamental/rudimentary component. 

        \item 
            However, blockchain is based on the idea of transactions. Everything
            is treated in terms of UTXOs, i.e. unspent transaction outputs. The
            current simulator introduces the idea of transactions and mempool
            (transactions that haven't yet been included in the block). 

        \item Introduction of Mempool enables us to implement and test the newer
            block propagation methods (COMPACT, XTHIN and GRAPHENE). 

        \item Orphan rates in the network determine whether it is profitable to
            mine blocks on a particular blockchain network.

        \item With addition of newer propagation methods on the simulator, we
            can now test bigger block sizes that original networks can't now.
            One of the examples of the current simulations being a 10MB standard
            block results in 90\% orphan rate but when run using a compact block,
            a 0-1\% orphan rate was noticed. These kind of scenarios can't be
            tested on the original main chain since block size changes require a
            hard fork and are not reversible.

        \item Higher orphan rates, i.e. 2-3\% can still be profitable for miners,
            which implies that the block size can be further adjusted to a much
            larger size which we are in the process of determining. 

        \item The above results are for dash blockchain that has a block interval
            for 2.5 minutes. In blockchain networks that have a higher block
            interval, even larger block sizes can be propagated without
            significant orphan rates. 

        \item Through this simulator, when a saturated block size is retrieved
            with acceptable orphan rate, a new chain can be forked for the most
            optimized results. 

    \end{enumerate}

    \subsubsection{The Emulator: } 
    \begin{enumerate}[label=(\alph*)]

        \item The future work for of the project is to build an emulator that
            will help implement off-chain state-of-the-art scaling methods and
            open ease of performing experiments for future ASU blockchain
            research projects.The emulator will be based on Hyper-V
            Kubernetes/Docker containers that will act as full-nodes, SPV and
            miners. The mining since it is compute intensive and provides
            deterministic results will be simulated by
            piecewise\_constant\_distribution by sampling the current data from
            Dash/Ethereum block explorers. 

        \item The prime focus of the emulator will be to test three techniques:
            bitcoin lightning network, Ethereum plasma and Sharding. 

        \item The lightning network is kicked off on the Bitcoin chain but only
            some group of users are using it. Plasma on the contrary is only
            theoretical in concept and Ethereum organization is working on
            funding organizations for the help. 

        \item The lightning network and plasma are both based on same grounds,
            and both claim billions of transactions as throughput.

        \item The aim of the emulator is to test 3 major hypothesis and 1 future
            work:
            \begin{enumerate}
                \item Unlimited transactions on the network don't necessarily
                    mean unlimited users. The network size bottleneck needs to
                    be calculated to figure out the number of users that can
                    comfirtably use the network without having high transaction
                    fee.

                \item In order to have billions of transactions, millions of
                    payment channels will be opened, which will bottleneck the
                    system. Finding the correlation between users, payment
                    channels and transactions that can be on the system.

                \item Although, billions of payments can happen on the lightning
                    network, all transactions are enforceable on the main
                    blockchain. Having a high throughput system base on a low
                    throughput system can lead to catastrophic scenarios. The
                    hypothesis that needs to be tested is for disaster recovery,
                    where if multiple payment channels close at once, if the
                    main chain will bottleneck based on the number of
                    transactions it will have to process. 

                \item Efficient routing between n-to-m endpoints means that the
                    most efficient (in terms of minimizing) the number of hops
                    configuration of the network would be large payment channel
                    hubs that would have open channels to everyone. Efficient
                    payment hubs lead to pseudo banks and existing banks would
                    be best suitable to run these payment hubs. Double spending
                    in lightning network is prevented by continuous monitoring
                    which is the current security model of the banks rather than
                    the original bitcoin network.
            \end{enumerate}

    \end{enumerate}


\section{Implementation Plan}
The simulator mentioned in the above component list is already implemented and
validated with all features in place. The simulator was first validated to the
real world parameters, which are: block interval and block size that affect
orphan rate and propagation time.

The emulator implementation schedule is listed below:    

\subsection{Proposed Schedule}
\begin{tabular}{|p{5cm}||p{5cm}||p{5cm}|}
    \hline
    \textbf{Task} & \textbf{Begin date} & \textbf{End date} \\
    \hline
    Setting up the Dash/Ethereum/Bitcoin system on Kubernetes & March 14th, 2018 & March 28th, 2018 \\
    \hline
    Testing scenarios for the above Blockchains & March 28th, 2018 & April 7th, 2018 \\
    \hline
    Building lightning network & April 7th, 2018 & April 28th, 2018 \\
    \hline
    Building Plasma & April 28th, 2018 & May 20th, 2018 \\
    \hline
    Testing and Results & May 21st,2018 & July 21st, 2018 \\
    \hline
\end{tabular}

\section{Evaluation Criteria}
The specific evaluation criteria for thesis should be based on the below
identified metrics:

\begin{enumerate}
    \item \textbf{Throughput: } The maximum throughput that can be achieved
        without having an orphan rate more than 0-2\%. Alternatively, the block
        size and orphan rate correlation can be calculated with respect to the
        total hash power of the network and the performance can be maximized. 

    \item \textbf{Scalability: } How many nodes can function on the bitcoin/dash
        network with sufficiently high block sizes to keep the throughput
        between 1-2\% variance from the maximum throughput. 
        What is the number of payment channels that can function without having
        bottlenecks or centralization. 

    % \item \textbf{Security Trade-offs:} What security trade offs does the
    %     proposed technique do, in order to reach this peak?
    % \item \textbf{Feature trade off:} What beneficial features does the proposed
    %     technique let go in order to reach the peak performance?
    % \item \textbf{Economic Incentive:} What orphan rate can be allowed on the
    %     network as a correlation of 

\end{enumerate}

\begin{thebibliography}{00}
    \bibitem{bitcoin} Nakamoto, Satoshi. "Bitcoin: A peer-to-peer electronic cash
        system." (2008).
    \bibitem{ethY} Wood, Gavin. "Ethereum: A secure decentralised generalised
        transaction ledger." Ethereum Project Yellow Paper 151 (2014): 1-32.
    \bibitem{ethW} Buterin, Vitalik. "Ethereum white paper." GitHub repository
        (2013).
    \bibitem{bgp} Lamport, L., Shostak, R., & Pease, M. (1982). The Byzantine
        generals problem. ACM Transactions on Programming Languages and Systems
        (TOPLAS), 4(3), 382-401.
    \bibitem{security} Gervais, Arthur, et al. "On the security and performance
        of proof of work blockchains." Proceedings of the 2016 ACM SIGSAC
        Conference on Computer and Communications Security. ACM, 2016.
    \bibitem{plasma} Poon, Joseph, and Vitalik Buterin. "Plasma: Scalable
        Autonomous Smart Contracts." White paper (2017).
    \bibitem{lightning} Poon, Joseph, and Thaddeus Dryja. "The bitcoin lightning
        network: Scalable off-chain instant payments." draft version 0.5 9
        (2016): 14.
    \bibitem{xthin} Peter Tschipper, BUIP010: Xtreme Thinblocks, 2016, Github
        Repository, https://github.com/BitcoinUnlimited/BUIP/blob/master/010.mediawiki  
    \bibitem{compact} Matt Corallo, Compact Block Relay, 2016, Github
        Repository, https://github.com/bitcoin/bips/blob/master/bip-0152.mediawiki  
    \bibitem{graphene} Ozisik, A. Pinar, et al. "Graphene: A New Protocol for
        Block Propagation Using Set Reconciliation." Data Privacy Management,
        Cryptocurrencies and Blockchain Technology. Springer, Cham, 2017.
        420-428.
    \bibitem{dash} Duffield, Evan, and Daniel Diaz. "Dash: A PrivacyCentric
        CryptoCurrency." (2014).
    \bibitem{segwit} Eric Lombrozo, Johnson Lau and Pieter Wuille, Segregated
        Witness (Consensus layer), 2015, Github Repository, https://github.com/bitcoin/bips/blob/master/bip-0141.mediawiki 
\end{thebibliography}

\end{document}
